% !TEX root = ../main.tex
\section{Implications for Evaluation}
\label{sec:evaluation}

The framework yields a concrete reorientation of evaluation practice, in four
directives and a checklist.

\subsection{Measure the grounding rate, and know where you can}
\label{sec:measure-rate}

Measure hallucination as \eqref{eq:rate}, via attribution audit. The procedure:
sample assertions from deployment-representative prompts; for each, test
condition (2) of Definition~\ref{def:grounding}---does the assertion trace to
identifiable support that is applicable at $c$, and does it respond appropriately
when that support is ablated, contradicted, or perturbed? The resulting statistic
separates the mechanism question, is the asserter evidence-coupled, from the
corpus question, is the evidence any good. Truth-matching conflates them, and
Corollary~\ref{cor:no-bound} shows the conflation cannot be repaired by
collecting more accuracy data.

The caveat of Section~\ref{sec:measurement} is a design constraint, not a
footnote. For retrieval-provided evidence the test is direct: ablate the context
and observe. For parametric encodings it is not, because the counterfactual ranges
over training corpora; attribution methods approximate it, imperfectly, and the
gap between the two cases is one of kind. So the honest position is that
\eqref{eq:rate} is estimable for retrieval-grounded assertion and largely not for
parametric assertion. Anyone reporting a hallucination rate should say which of
the two they measured.

A note on adjacent instruments. Semantic-entropy methods detect a related defect
by clustering sampled outputs into meaning classes via bidirectional entailment
\citep{farquhar2024}. That is a useful signal, and it is not this measurement:
clustering by meaning deliberately quotients out the formulation-level variation
that Section~\ref{sec:re-failure} identifies as constitutive. Semantic entropy
measures instability of content under resampling; \eqref{eq:rate} measures
decoupling of assertion from evidence. A system can be low-entropy and
ungrounded---confidently and consistently wrong---so the two belong side by side
rather than one standing in for the other.

\subsection{Score the off-diagonal cells correctly}
\label{sec:score-cells}

Cell (iii) outputs---grounded falsehoods---should be charged to the evidential
base rather than the assertion mechanism. They call for corpus and retrieval
intervention and should not depress a hallucination metric. Cell (ii)
outputs---veridical hallucinations---should be charged to the mechanism despite
their truth. A benchmark that scores them as successes is training systems to
guess confidently.

Operationally, and this is the most directly implementable consequence of the
paper: \textbf{evaluation suites should include items whose correct answers are
absent from the accessible evidence, where the only correct behaviour is
abstention, and should penalise true answers on those items.} A system that
answers such an item correctly has demonstrated nothing except that its guessing
distribution happened to align with the world, and rewarding it selects for
exactly the decoupling of assertion from evidence that
Section~\ref{sec:unity} identifies as the root pathology.

The same measurability asymmetry applies here and is easy to miss. Establishing
that an answer is absent from the accessible evidence is straightforward when you
control the retrieval corpus, and not generally possible for a frontier model's
parametric knowledge. Abstention items are therefore constructible for
retrieval-grounded systems and, for unaugmented models, only in domains where
corpus absence can be independently established. This is a limitation of the
method, and another reason the architecture of
Section~\ref{sec:warrant} is not merely good hygiene.

\subsection{Audit congruentiality directly}
\label{sec:re-audit}

Since RE-failure is the root pathology, paraphrase-equivariance deserves promotion
from robustness footnote to headline metric. For a battery of logically equivalent
formulation pairs---contraposition, negation duality, notation shifts, active
against passive---measure divergence in assertion behaviour at matched output
probability. The statistic is in effect an empirical estimate of how far $\Out_M$
sits from classicality: a scalar summary of the distance between the system's
operator and a proposition-directed one.

Two methodological cautions. Matching output probability across paraphrases is
itself non-trivial, since the probability is formulation-sensitive too. And the
battery must control for the possibility that paraphrases differ in tokenisation
difficulty rather than in logical form, or the statistic will measure the wrong
thing.

Applied across inference-compute budgets and training checkpoints, the same
battery answers the question left open in Section~\ref{sec:reasoning}: whether
scale, post-training, and extended deliberation move systems toward normality or
merely broaden the set of formulations handled. Falling divergence is convergence
on classicality. Steady divergence under rising accuracy is coverage. I know of no
systematic measurement of this, and it is the experiment the framework most wants
run.

\subsection{Reformulate subject-level appraisal}
\label{sec:appraisal}

Metrics and documentation phrased in terms of what a model ``knows'' or
``believes'' should be restated in the $(\Out_M, \Grd_M)$ idiom: assertion rates,
grounding rates, stability profiles. Section~\ref{sec:deflationism} gives the
reason---the doxastic idiom carries normality assumptions that distort measurement
design and downstream reliance.

I would not extend this to every evaluative term with a mentalistic ring. Notions
like deception and honesty have operational definitions in safety evaluation that
do real work, and the framework does not impugn them; it reconstructs them. The
safety-relevant phenomenon is precisely the dissociation of
Section~\ref{sec:dissociation}: divergence between what a model's internal state
registers and what it asserts. That is expressible in the present idiom without
loss, and expressing it that way makes clear why it is a defect---not insincerity
by a subject, but assertion unaccountable to the system's own best estimate.

\subsection{What an architecture would have to achieve}
\label{sec:checklist}

The framework's most useful output may be constructive. It specifies what a
system would have to do to earn doxastic description:

\begin{enumerate}
  \item maintain content-level rather than formulation-level representations,
        sufficient to validate RE;
  \item implement processes whose \emph{function} is closure and consistency
        maintenance, delivering K, M, and D at the level of competence;
  \item support self-access channels causally coupled to first-order states,
        delivering 4 and 5; and
  \item embed assertion in a persistent cross-context economy that makes closure
        a functional norm rather than an ornament---which, by
        Section~\ref{sec:human-parallel}, is what makes (1)--(3) norms rather than
        accidents.
\end{enumerate}

Retrieval-grounded generation, verifier-loop architectures, outcome-reinforcement
training, and persistent-memory agents are partial steps along distinct axes, and
the framework says which axis each addresses. Retrieval supplies an explicit $E$
and makes \eqref{eq:rate} measurable; it does nothing for congruentiality.
Verifier loops address (2) and can be checked against
Section~\ref{sec:re-audit}. Persistent memory is a precondition of (4) and not
sufficient for it, since storage is not recruitment. The apparatus turns ``do
these systems believe?'' from a slogan-war into a checklist, and the checklist has
the useful property that each item names its own measurement.
